\documentclass[11pt]{article}
\usepackage{amsmath,amssymb}
\usepackage{geometry}
\usepackage{setspace}
\usepackage{enumitem}
\geometry{margin=1in}
\setstretch{1.15}

\begin{document}

\begin{center}
{\Large \textbf{CSE 400 — Lecture 23 Scribe}}\\
\vspace{0.2em}
Predicting the Impossible — Tail Bounds and System Reliability
\end{center}

\section{Problem Context: System Reliability Under Uncertainty}

\subsection{Scenario Definition}

The system models Amazon Black Friday Sale 2026 traffic.

High-volume distributed systems must handle random request arrivals.

Requests arrive at servers as random processes.

Modeled using Poisson arrivals:
\[
X \sim \text{Poisson}(\lambda)
\]

Critical system event:

\[
\text{Server overload occurs if } X \ge k
\]

where:
\begin{itemize}[noitemsep]
\item $X$ : number of requests in a time window
\item $k$ : system crash threshold
\end{itemize}

\section{Tail Probability as Failure Metric}

\subsection{Definition}

\[
P(X \ge k)
\]

The system reliability question reduces to computing this probability.

Key statement:

\[
\text{“The mean is safe. The tail is catastrophic.”}
\]

\subsection{Example from Simulation}

\[
\lambda = 50,\quad k = 60
\]

Empirical crash rate:
\[
8.7\% \quad (\text{87 crashes in 1000 minutes})
\]

Thus:
\[
\text{Goal: Estimate } P(X \ge 60)
\]

\section{Necessity of Probability Bounds}

Exact computation of $P(X \ge k)$ may not always be feasible.

Thus, upper bounds are used.

\subsection*{Hierarchy of Bounds}

\begin{center}
\begin{tabular}{c|c|c}
Level & Method & Required Information \\ \hline
1 & Markov & Mean \\
2 & Chebyshev & Mean, Variance \\
3 & Chernoff & Full distribution (MGF)
\end{tabular}
\end{center}

Observation:
\[
\text{More information} \Rightarrow \text{tighter bounds}
\]

\section{Markov’s Inequality}

\subsection{Statement}

For non-negative random variable $X$:

\[
P(X \ge a) \le \frac{E[X]}{a}
\]

\subsection{Inputs}

\[
\mu = E[X]
\]

\subsection{Behavior}

Decay: linear in $1/a$

Example:
\[
\text{Exact} = 0.004,\quad \text{Markov} = 0.714
\]

\subsection{Interpretation}

\begin{itemize}[noitemsep]
\item Bound is loose
\item Strong overestimation
\item Leads to pessimistic system design
\end{itemize}

\section{Chebyshev’s Inequality}

\subsection{Statement}

\[
P(|X - \mu| \ge a) \le \frac{\sigma^2}{a^2}
\]

Equivalent tail form:

\[
P(X \ge a) \le \frac{\sigma^2}{(a - \mu)^2}
\]

\subsection{Inputs}

\[
\mu = E[X], \quad \sigma^2 = \text{Var}(X)
\]

\subsection{Behavior}

Decay: polynomial in $1/a^2$

Example:
\[
\text{Exact} = 0.004,\quad \text{Chebyshev} = 0.125
\]

\subsection{Interpretation}

\begin{itemize}[noitemsep]
\item Tighter than Markov
\item Still significantly overestimates
\end{itemize}

\section{Chernoff Bound}

\subsection{Statement (General Form)}

\[
P(X \ge a) \le \min_{t>0} e^{-ta} M_X(t)
\]

where:
\[
M_X(t) = E[e^{tX}]
\]

\subsection{Inputs}

Full distribution via moment generating function.

\subsection{Behavior}

Decay: exponential

Example:
\[
\text{Exact} = 0.004,\quad \text{Chernoff} = 0.092
\]

\subsection{Interpretation}

\begin{itemize}[noitemsep]
\item Tightest bound among the three
\item Closely follows actual tail probability
\end{itemize}

\section{Comparative Decay Rates}

\begin{center}
\begin{tabular}{c|c|c}
Method & Decay Type & Expression \\ \hline
Markov & Linear & $1/k$ \\
Chebyshev & Polynomial & $1/k^2$ \\
Chernoff & Exponential & $e^{-k}$
\end{tabular}
\end{center}

\[
\text{Chernoff} \ll \text{Chebyshev} \ll \text{Markov}
\]

\section{Case Study: SLA Breach}

\subsection{Parameters}

\[
\lambda = 500,\quad k = 600
\]

\subsection{Objective}

\[
\text{Estimate } P(X \ge 600)
\]

Exact:
\[
1.4\%
\]

\section{Cost of Loose Bounds}

\begin{center}
\begin{tabular}{c|c}
Method & Estimated Failure Probability \\ \hline
Chernoff & 8.8\% \\
Chebyshev & 20\% \\
Markov & 90.9\%
\end{tabular}
\end{center}

Observation:

Markov exaggerates risk drastically.

\section{Financial Interpretation}

\subsection{Waste Multiplier}

\[
\text{Waste Multiplier} =
\frac{\text{Bound Estimate}}{\text{Exact Probability}}
\]

\subsection{Example at $\lambda = 500$}

\begin{itemize}[noitemsep]
\item Chernoff: $\approx 10.7\times$
\item Chebyshev: $\approx 6422\times$
\item Markov: $\approx 107039\times$
\end{itemize}

Conclusion:

\[
\text{Loose bound} \Rightarrow \text{Over-provisioning} \Rightarrow \text{Financial loss}
\]

\section{Distribution Sensitivity}

Observation:

Chernoff performs well for:
\begin{itemize}[noitemsep]
\item Poisson
\item Binomial
\end{itemize}

For heavy-tailed distributions (e.g., Gamma):

Chernoff may overestimate.

\section{Engineer’s Decision Matrix}

\begin{center}
\begin{tabular}{c|c|c|c}
Method & Requires & Decay & Use Case \\ \hline
Markov & Mean & Linear & Only mean known \\
Chebyshev & Mean, Variance & Polynomial & Variance known \\
Chernoff & Full MGF & Exponential & Precise SLA sizing
\end{tabular}
\end{center}

\section{Class Exercise: Laplace Distribution}

\subsection{Given}

\[
X \sim \text{Laplace}(\mu, b), \quad \sigma = b\sqrt{2}
\]

\subsection{Objective}

\[
a = \mu + 3\sigma
\]

Compute Chernoff bound for:
\[
P(X \ge a)
\]

\subsection{Derivation Blueprint}

\begin{enumerate}[noitemsep]
\item Compute MGF:
\[
M_X(t) = E[e^{tX}]
\]
\item Apply Chernoff:
\[
P(X \ge a) \le e^{-ta} M_X(t)
\]
\item Optimize:
\begin{itemize}[noitemsep]
\item Differentiate w.r.t $t$
\item Set derivative to $0$
\item Solve for $t^*$
\end{itemize}
\item Substitute optimal $t^*$
\end{enumerate}

\section{Final Synthesis}

\begin{enumerate}[noitemsep]
\item System failures are governed by tail probabilities:
\[
P(X \ge k)
\]
\item Exact computation often infeasible $\Rightarrow$ bounds required.
\item Trade-off:
\[
\text{Less information} \Rightarrow \text{weaker bound} \Rightarrow \text{higher cost}
\]
\[
\text{More information} \Rightarrow \text{tighter bound} \Rightarrow \text{efficient design}
\]
\item Mathematical decay governs engineering decisions:
\[
\text{Linear} \rightarrow \text{Polynomial} \rightarrow \text{Exponential}
\]
\item Core principle:
\[
P(X \ge a) \le e^{-t^* a} M_X(t^*)
\]
\end{enumerate}

\end{document}